\documentclass[10pt,a4paper]{amsart}
\usepackage[margin=19mm]{geometry}
\usepackage{amsmath,amssymb,mathtools}
\usepackage[scr=rsfs,cal=euler]{mathalfa}
\usepackage{xfrac}
\usepackage[unicode,hidelinks,bookmarks=false]{hyperref}
\newcommand{\ie}{i.e.\ }
\newcommand{\cf}{cf.\ }
\newcommand{\resp}{respectively\ }
\newcommand{\Subsection}{\subsection}
\providecommand{\nobreakdash}{\nobreak\hbox{-}}
\setlength{\emergencystretch}{2em}
\multlinegap0pt
\usepackage{amssymb}
\usepackage{enumerate}
\newtheorem{theorem}{Theorem}
\def\bM{\mathbb{M}}
\newcommand{\qsym}{\mathrm{qsym}}
\newcommand{\sym}{\mathrm{sym}}
\title{Averages over matrix unitary orbits \\ and spectral order}
\author{Jean-Christophe Bourin, Eun-Young Lee}
\date{Source: arXiv:2606.15624v2 (CC BY 4.0)}
\begin{document}
\maketitle

\noindent\emph{Source and licence.} Averages over matrix unitary orbits \\ and spectral order. Version arXiv:2606.15624v2 (CC BY 4.0), distributed by arXiv under the Creative Commons Attribution 4.0 International licence, http://creativecommons.org/licenses/by/4.0/, as recorded in the arXiv licence metadata for this exact version and shown on the abstract page https://arxiv.org/abs/2606.15624v2. Full licence notice: \texttt{licenses/arxiv-2606.15624-CC-BY-4.0.txt}.\par\smallskip

\noindent\textit{Editorial scope.} Counted unit: the article's theorem th2 (source0.tex lines 249--253). The article's complete original proof of it is delivered with it (source0.tex lines 394--407, one \texttt{proof} environment of 14 lines, which the article places away from the statement). No mathematics is written, repaired or re-derived: the statement and the proof are the article's own lines, in the article's own notation, and no retained line is substituted or re-flowed.  No further source-local context is needed: every cross-reference of every retained line resolves inside the two delivered spans.  Nothing was omitted from any retained span, so the package carries no omission record.  No retained line contains a citation, so the wrapper carries no bibliography and no bibliography entry.  No retained line contains a figure environment, an image-inclusion command or drawing code, so no image is delivered and the package declares no asset dependency.  Editorial formatting only, in addition: an AMS \texttt{amsart} preamble that restates the article's own declarations and macros used by the retained lines, namely the environments \texttt{theorem} on separate counters, together with the standard packages those lines need.  The retained environments print in this document as Theorem 1 in order of appearance, whereas the article prints the same items with the numbering of its own full document.  The complete native source of the article as distributed in the arXiv e-print for this exact version is bundled unchanged as \texttt{sources/202878/source0.tex}.
\begin{theorem}\label{th2} Let $f,g:[0,\infty)\to[0,\infty)$ be nondecreasing. Assume that $f(t)$ is concave and  $g(t)$ is convex. If $\{A_k\}_{k=1}^m$ are  normal matrices in $\bM_d$, then there exist unitary matrices $\{V_i\}_{i=1}^d$  in $\bM_d$ such that 
$$ 
g\circ f\left(\left|\sum_{k=1}^m A_k\right|\right) \le \frac{1}{d} \sum_{i=1}^d V_i g\left(\sum_{k=1}^m f(|A_k|)\right)V_i^*.
$$
\end{theorem}

\begin{proof} To see that $\sqrt{2}$ is optimal, it suffices to pick
$$
Z=\begin{pmatrix} 0&2 \\ 0&0\end{pmatrix}
$$
as
$$
|Z|_{\qsym} =  \begin{pmatrix} \sqrt{2}&0 \\ 0&\sqrt{2}\end{pmatrix}, \qquad |Z|_{\sym} =  \begin{pmatrix} 1&0 \\ 0&1\end{pmatrix}.
$$
To establish the main conclusion, apply Theorem \ref{th2} with two matrices $A_1=|Z|^2$, $A_2=|Z^*|^2$, and the functions $f(t)=\sqrt{t}$ and $g(t)=t$. This gives
$$
\sqrt{|Z|^2+|Z^*|^2} \le \frac{ 1}{d}  \sum_{i=1}^d V_i \left(|Z|+|Z^*|\right)V_i^*
$$
and dividing both sides by $\sqrt{2}$ yields the result.
\end{proof}

\end{document}
